\section{Conclusion and Future Work}

\subsection{Summary of Contributions}

This project demonstrates that a carefully designed, lightweight
convolutional neural network can achieve state-competitive crop disease
classification accuracy while enabling real-time inference on resource-
constrained edge devices. The key contributions are:

\begin{enumerate}[label=\textbf{C\arabic*}, leftmargin=2.5em, itemsep=4pt]
  \item Designed and validated a custom 6.3M-parameter CNN achieving 97.8\%
        accuracy on 18-class crop disease classification — within 0.4 pp of
        ResNet-50 at 4$\times$ fewer parameters.
  \item Established a rigorous benchmarking protocol comparing custom,
        transfer-learned, and lightweight architectures on identical data
        splits, preprocessing, and evaluation metrics.
  \item Demonstrated that data augmentation and batch normalisation
        contribute the largest individual gains (+3.7 pp and +4.3 pp
        respectively) in the PlantVillage domain.
  \item Deployed the model as a Flask-based web application with TensorFlow
        Serving, enabling smartphone-based disease diagnosis with measured
        end-to-end latency under 300\,ms (server round-trip).
  \item Open-sourced all code, trained weights, and documentation to
        facilitate reproducibility and extension by the research
        community.\footnote{Repository available at
        \url{https://github.com/crop-disease-detection/cdd}}
\end{enumerate}

\subsection{Limitations}

Several limitations of the current work merit acknowledgment. First, the
PlantVillage dataset consists of laboratory images with uniform backgrounds;
performance degrades on in-field images with complex canopies, shadows, and
soil backgrounds. Second, the system classifies a single leaf at a time and
does not detect multiple diseases co-occurring on the same plant — a common
real-world scenario. Third, the web application requires server-side
inference, limiting utility in areas with poor internet connectivity. Fully
on-device deployment (TensorFlow Lite) was explored but not productionised.

\subsection{Future Work}

Future work will pursue five directions:

\begin{enumerate}[label=\textcolor{accent}{\textbullet}, leftmargin=2em, itemsep=3pt]
  \item \textbf{Field Dataset Collection:} Partner with agricultural
        extension services to curate a large-scale, in-field image dataset
        with diverse backgrounds, lighting conditions, and disease severity
        levels.
  \item \textbf{On-Device Deployment:} Convert the custom CNN to TensorFlow
        Lite with 8-bit quantisation and benchmark accuracy-latency trade-offs
        on mid-range Android devices.
  \item \textbf{Multi-Disease Detection:} Extend the architecture to object
        detection (YOLO-based or Faster R-CNN) to localise multiple disease
        lesions within a single image.
  \item \textbf{Severity Estimation:} Augment classification with regression
        heads to estimate disease severity (percentage leaf area affected),
        enabling targeted treatment recommendations.
  \item \textbf{Multilingual Farmer Interface:} Develop localised mobile
        interfaces with voice-based interaction in regional languages to
        maximise accessibility for non-literate farmers.
\end{enumerate}

% ══════════════════════════════════════════════════════════════════════════
% ── References ────────────────────────────────────────────────────────────
% ══════════════════════════════════════════════════════════════════════════
\newpage
\begin{thebibliography}{99}

\bibitem{mohanty2016}
S.\ P.\ Mohanty, D.\ P.\ Hughes, and M.\ Salathé,
``Using deep learning for image-based plant disease detection,''
\textit{Frontiers in Plant Science}, vol.\ 7, p.\ 1419, 2016.

\bibitem{sladojevic2016}
S.\ Sladojevic, M.\ Arsenovic, A.\ Anderla, D.\ Culibrk, and D.\ Stefanovic,
``Deep neural networks based recognition of plant diseases by leaf image
classification,''
\textit{Computational Intelligence and Neuroscience}, vol.\ 2016, 2016.

\bibitem{ferentinos2018}
K.\ P.\ Ferentinos,
``Deep learning models for plant disease detection and diagnosis,''
\textit{Computers and Electronics in Agriculture}, vol.\ 145, pp.\ 311--318,
2018.

\bibitem{too2019}
E.\ C.\ Too, L.\ Yujian, S.\ Njuki, and L.\ Yingchun,
``A comparative study of fine-tuning deep learning models for plant disease
identification,''
\textit{Computers and Electronics in Agriculture}, vol.\ 161, pp.\ 272--279,
2019.

\bibitem{brahimi2017}
M.\ Brahimi, K.\ Boukhalfa, and A.\ Moussaoui,
``Deep learning for tomato diseases: Classification and symptoms
visualization,''
\textit{Applied Artificial Intelligence}, vol.\ 31, no.\ 4, pp.\ 299--315,
2017.

\bibitem{kamilaris2018}
A.\ Kamilaris and F.\ X.\ Prenafeta-Boldú,
``Deep learning in agriculture: A survey,''
\textit{Computers and Electronics in Agriculture}, vol.\ 147, pp.\ 70--90,
2018.

\bibitem{barbedo2018}
J.\ G.\ A.\ Barbedo,
``Factors influencing the use of deep learning for plant disease
recognition,''
\textit{Biosystems Engineering}, vol.\ 172, pp.\ 84--91, 2018.

\bibitem{ramcharan2017}
A.\ Ramcharan, K.\ Baranowski, P.\ McCloskey, B.\ Ahmed, J.\ Legg, and
D.\ P.\ Hughes,
``Deep learning for image-based cassava disease detection,''
\textit{Frontiers in Plant Science}, vol.\ 8, p.\ 1852, 2017.

\bibitem{fuentes2017}
A.\ Fuentes, S.\ Yoon, S.\ C.\ Kim, and D.\ S.\ Park,
``A robust deep-learning-based detector for real-time tomato plant diseases
and pests recognition,''
\textit{Sensors}, vol.\ 17, no.\ 9, p.\ 2022, 2017.

\bibitem{saleem2019}
M.\ H.\ Saleem, J.\ Potgieter, and K.\ M.\ Arif,
``Plant disease detection and classification by deep learning,''
\textit{Plants}, vol.\ 8, no.\ 11, p.\ 468, 2019.

\bibitem{thenmozhi2019}
K.\ Thenmozhi and U.\ S.\ Reddy,
``Crop pest classification based on deep convolutional neural network and
transfer learning,''
\textit{Computers and Electronics in Agriculture}, vol.\ 164, p.\ 104906,
2019.

\bibitem{lu2017}
Y.\ Lu, S.\ Yi, N.\ Zeng, Y.\ Liu, and Y.\ Zhang,
``Identification of rice diseases using deep convolutional neural networks,''
\textit{Neurocomputing}, vol.\ 267, pp.\ 378--384, 2017.

\bibitem{hughes2015}
D.\ P.\ Hughes and M.\ Salathé,
``An open access repository of images on plant health to enable the
development of mobile disease diagnostics,''
\textit{arXiv preprint arXiv:1511.08060}, 2015.

\bibitem{pydipati2006}
R.\ Pydipati, T.\ F.\ Burks, and W.\ S.\ Lee,
``Identification of citrus disease using color texture features and
discriminant analysis,''
\textit{Computers and Electronics in Agriculture}, vol.\ 52, no.\ 1--2,
pp.\ 49--59, 2006.

\end{thebibliography}
